\def\stat{voronov}



\def\tit{ДИНАМИЧЕСКОЕ СОЗДАНИЕ КАРТ УРОВНЯ WiFi-СИГНАЛОВ ДЛЯ СИСТЕМ ЛОКАЛЬНОГО

ПОЗИЦИОНИРОВАНИЯ$^*$}



\def\titkol{Динамическое создание карт уровня WiFi-сигналов

%для систем локального позиционирования

}



\def\autkol{Р.\,В. Воронов, С.\,В.~Малодушев}



\def\aut{Р.\,В. Воронов$^1$, С.\,В.~Малодушев$^2$}



\titel{\tit}{\aut}{\autkol}{\titkol}



{\renewcommand{\thefootnote}{\fnsymbol{footnote}}

\footnotetext[1]{Исследования проводились в рамках деятельности МИП ООО <<Наносети>> и ЗАО

<<РТЛ-Сер\-вис>>, а также поддерживались Петрозаводским государственным

университетом (Программа стратегического развития ПетрГУ на 2012--2016~гг.),

Министерством образования и науки РФ, Фондом СР МФП в НТС и Министерством

экономического развития Республики Карелии (Программа <<Разработка системы слежения

за мобильными абонентами распределенной WiFi-сети>> на 2013~г.), а также в рамках

программы УМНИК <<Анализ и проектирование WiFi-сетей в помещениях>> на

2012--2013~гг.}}



\renewcommand{\thefootnote}{\arabic{footnote}}

\footnotetext[1]{Петрозаводский государственный университет, rvoronov@sampo.ru}

\footnotetext[2]{Петрозаводский государственный университет, svm@petrsu.ru}



 \label{st\stat}



% \vspace*{-12pt}



\Abst{Карты уровня сигналов используются для локации беспроводных мобильных

устройств. Увеличение беспроводного сегмента сети требует создания новых карт уровня

сигналов для осуществления локации в новых областях. Традиционно подобные карты

формируются из большого количества измерений в определенных точках, подобный процесс

занимает много времени. В~статье исследована задача динамического создания карт уровня

сигналов. Описан эвристический алгоритм динамического создания карт уровня сигналов во

время эксплуатации системы локации. При оценивании средних значений уровня сигналов

использована функция, являющаяся аналогом функции правдоподобия. Разработанный

алгоритм позволяет автоматизировать процесс создания подобных карт. В~качестве

беспроводных сетей использовались WiFi-cе\-ти, работающие в частотном диапазоне

2,4~ГГц. Результаты тестов показали, что предложенный алгоритм может быть использован

для решения данной задачи.}



\KW{WiFi; базовые станции; опорная точка; уровень сигнала; локация; карта уровня

сигналов}



\DOI{10.14357\!/08696527140105}



\vskip 12pt plus 9pt minus 6pt



 \thispagestyle{headings}



\section{Введение}



 Как известно, спутниковые системы навигации невозможно использовать в

закрытых помещениях~--- в зданиях, шахтах, подземных парковках, переходах

метро и~т.\,д. Для таких случаев разрабатываются и применяются специальные

системы локации, основанные, например, на обработке уровня принимаемого

сигнала от стационарных точек доступа (базовых станций) беспроводной сети.

Предварительно помещение разбивается на множество опорных точек, для

каж\-дой из которых определяются карты уровня принимаемых сигналов от

базовых станций~[1--3].



 Поскольку уровень принимаемого сигнала сильно зависит от числа

отражений, любые изменения в помещении или смена ориентации антенн

базовых станций приводят к необходимости заново настраивать систему

локации~--- пересчитывать карты уровня сигналов для нескольких или всех

опорных точек помещения. Это трудоемкая операция, так как она требует

непосредственного обхода оператором системы опорных точек и определения

для каждой из них среднего значения уровня сигнала. В~рассматриваемой

работе предполагается автоматическое формирование карт уровня

принимаемых сигналов в ходе эксплуатации системы, что является

существенной особенностью рассматриваемого подхода и отличает его от

других подходов к решению задачи локации внутри помещений.



 Для систем локации, основанных на WiFi-сетях, существуют ограничения: с

одной стороны, в каждой опорной точке помещения нужно слышать

достаточное для локации число базовых станций (желательно больше трех), с

другой стороны, используемые базовыми станциями диапазоны частот по

возможности не должны пересекаться, иначе количество интерферирующих

сигналов будет высоко.



 Введем необходимые обозначения. Пусть $N$~--- число опорных точек

помещения, в которых может находиться мобильный объект, $D$~--- число

базовых станций. На приемной стороне уровни сигналов от базовых станций

являются дискретными величинами, поскольку радиопередача ведется в

цифровой форме. Тогда в ходе настройки системы необходимо для каждой

базовой станции~$d$ оценить распределение $p_d(x\vert j)$ условных

вероятностей уровней сигнала~$x$ в опорных точках $j\hm= 1,\ldots , N$.



 В ходе эксплуатации системы местоположение мобильного объекта

определяется следующим образом. Пусть $p_j$~--- априорная вероятность

нахождения объекта в опорной точке~$j$ ($j\hm= 1,\ldots , N$). Предположим,

что зарегистрирован вектор уровней сигнала $\overline{x}\hm= (x_1,\ldots

,x_d)$, где $x_d$~--- уровень сигнала от базовой станции~$d$. Предполагая

независимость уровней сигнала от различных базовых станций, апостериорную

вероятность~$q_j$ нахождения объекта в опорной точке~$j$ можно найти по

формуле:

 \begin{equation}

 q_j= \fr{p_j\prod\limits_{d=1}^D p_d(x_d\vert j)} {\sum\limits_{k=1}^N p_k

\prod\limits_{d=1}^D p_d(x_d\vert k)}\,.

 \label{e1-vor}

 \end{equation}



 Как правило, в качестве искомой опорной точки выбирается та, для которой

апостериорная вероятность принимает наибольшее значение.



 В данной статье предлагается эвристический метод, позволяющий без

перенастройки оператором системы локации в ходе ее эксплуатации

автоматически оценивать вероятности $p_d(x\vert j)$ для добавляемых в

систему новых базовых станций, используя значения~$q_j$, полученные на

основе имеющихся (<<старых>>) базовых станций. Предложенный метод

опробован в системе локации, основанной на WiFi-сетях.



\section{Построение областей локации }



 Рассмотрим существующие частотные ограничения, накладываемые на

количество используемых базовых станций. Сети WiFi работают в частотном

диапазоне 2400--2495~МГц. Данный частотный диапазон делится на

14~каналов, несущая частота каждого из которых отличается на 5~МГц, при

этом каждый радиоканал имеет фиксированный частотный диапазон шириной

22~МГц. Для того чтобы сигналы от базовых станций не интерферировали друг

с другом, при построении WiFi-се\-тей их несущие частоты должны отличаться

не менее чем на 25~МГц, т.\,е.\ в доступном для WiFi частотном диапазоне

2,4~ГГц возможна работа без взаимных помех лишь на трех

 радиоканалах~[4, 5]. В~данной статье не рассматривается случай

использования диапазона 5~ГГц.



 На рис.~1 представлен типичный пример деления WiFi-сети на

радиозоны~(\textit{а}) и пример тех же радиозон, используемых при

локации~(\textit{б}),~--- каждой базовой станции соответствует свой номер

канала, определяющий частотный диапазон, и радиозона. Расчет позиции

мобильного объекта возможен только в областях пересечения радиозон,

являющихся областями локации. В~данном случае областям локации

соответствуют радиозоны с номерами каналов~6 и~11 (см.\ рис.~1,\,\textit{б}),

несущие частоты данных радиозон не интерферируют друг с другом. Из

рис.~1,\,\textit{б} видно, что для систем локации плотность размещения

базовых станций выше в сравнении с традиционными WiFi-се\-тя\-ми, ввиду

необходимости организации областей локации.



 \begin{figure}[b] %fig1

 \vspace*{-3pt}

 \begin{center}

 \mbox{%

 \epsfxsize=59.4mm

 \epsfbox{vor-1.eps}

 }

 \end{center}

 \vspace*{-13pt}

 \Caption{Использование WiFi-каналов для организации радиозон}

 \end{figure}



 По мере удаления от источника излучения базовой станции радиосигнал

затухает, ввиду чего границы радиозон определяются чувствительностью

радиоприемника на стороне клиента.



 \section{Применение карт уровня сигналов}



 Карты уровня сигналов используются в расчетах при локации, и их

формирование является необходимой и ресурсоемкой задачей. Рассмотрим, как

осуществляется создание и применение данных карт при расчетах в локации.



 В момент времени~$t$ каждой опорной точке~$j$ соответствует уникальный

вектор уровней сигналов $\overline{x}\hm= (x_1,\ldots , x_d)$ для радиозон

соответствующих базовых станций. Для каждой базовой станции формируется

карта, содержащая усредненные показания принятого от нее сигнала в

заданном наборе опорных точек $j\hm= 1,\ldots , N$. 



При локации

осуществляется сравнение полученного в конкретной точке вектора уровней

сигналов RSSI (Received Signal Strength Indication) со значениями, хранимыми

во всех картах. По результатам данного сравнения определяется наиболее

вероятная позиция мобильного объекта~[6].





 \begin{floatingfigure}[t]{65mm} %fig2

 \vspace*{1pt}

 \begin{center}

 \mbox{%

 \epsfxsize=59.898mm 

 \epsfbox{vor-2.eps}

 }

 \end{center}

 \vspace*{-6pt}

\Caption{Положение мобильного объекта относительно базовых станций}

\vspace*{9pt}

\end{floatingfigure}



 Для корректного расчета позиции мобильного объекта необходимо, чтобы

данный объект принимал пока-\linebreak зания RSSI не менее чем от трех\linebreak базовых

станций, для которых сформированы карты уровня сигналов, при этом данный

мобильный объект должен находиться внутри области, по периметру которой

располагаются эти базовые станции (рис.~2).



 На рис.~2 представлены две опорные точки~--- $A$ и~$B$. Локация

мобильного объекта~$D$ в опорной точке~$A$ осуществляется на основе

информации об уровнях сигналов от базовых станций $A1$, $A2$ и~$A3$. 



При

локации в опорной точке~$A$ также используются показания уровней сигналов

от базовой станции~$A4$ в случаях, когда мобильный объект~$D$ принимает

радиосигнал от данной базовой станции, при этом, ввиду частотных

ограничений, базовые станции~$A1$ и~$A4$ по возможности не должны

находиться в зоне слышимости друг друга, так как их сигналы будут

интерферировать друг с другом.





 Рассматриваемая модель позволяет формировать карты уровня сигналов в

ходе эксплуатации системы с целью их дальнейшего использования при

локации, что делает возможным автоматизировать процесс создания подобных карт и,

как следствие, увеличить количество областей локации.



\section{Математическая модель}



 Пусть имеется некоторое число базовых станций, для которых известны

распределения вероятностей уровней сигнала в каждой опорной точке

помещения. Тогда, обрабатывая по формуле~(\ref{e1-vor}) для мобильного

объекта уровни сигналов от базовых станций, можно определять вероятности

расположения объекта в опорных точках помещения. Пусть добавляется новая

базовая станция, для которой не определены распределения вероятностей

уровней сигналов в опорных точках.



 Предположим для простоты, что уровни сигнала от новой базовой станции

являются непрерывными случайными величинами $\xi_1, \ldots , \xi_M$,

которые в результате $M$ независимых измерений приняли значения

$x_1,\ldots , x_M$.



 В модели примем плотности распределений уровней сигналов в опорных

точках~$j$ равными



\noindent

 \begin{equation}

 p(x\vert j; a_j) =\fr{1}{\sigma\sqrt{2\pi}}\exp\left( - \fr{(x-

a_j)^2}{2\sigma^2}\right)\,, \enskip j=1,\ldots , N\,,

 \label{e2-vor}

 \end{equation}

где параметр $\sigma$ является известным и одинаковым для всех опорных

точек~$j$. Теоретически в этом случае допускается ненулевая вероятность для

отрицательных значений уровня сигналов, но она является ничтожно малой.

Параметры $a_j$~--- математические ожидания значений уровня сигнала в

опорных точках $j\hm= 1,\ldots , N$~--- априори не известны. Задача состоит в

поиске их точечных оценок.



 Пусть $q_{ij}$~--- известные вероятности нахождения объекта в опорной

точке~$j$ в момент регистрации $i$-го измерения, определенные по

формуле~(\ref{e1-vor}).



 Плотность распределения уровня сигнала~$x$ в момент регистрации $i$-го

измерения вычисляется по формуле полной вероятности:

 $$

 p_i(x) =\sum\limits_{j=1}^N p(x\vert j;a_j)q_{ij}\,.

 $$



 Составим функцию, являющуюся аналогом логарифмической функции

правдоподобия для измерений $x_1,\ldots , x_M$:

 \begin{equation}

 L(\overline{x};\overline{a}) =\sum\limits_{i=1}^M \ln \left( \sum\limits_{j=1}^N

p(x_i\vert j;a_j)q_{ij} \right)\,,

 \label{e3-vor}

 \end{equation}

где $\overline{x}=(x_1,\ldots , x_M)$ и $\overline{a}\hm= (a_1,\ldots , a_N)$.

В~случае, если вероятности~$q_{ij}$ меняются от измерения к измерению,

случайные величины $\xi_1,\ldots , \xi_M$ будут распределены по-раз\-но\-му, и

в строгом смысле функция~(\ref{e3-vor}) не является функцией правдоподобия.

Задача оценивания параметров в такой постановке сама по себе является

нетривиальной. Предлагается оценивать вектор~$\overline{a}$ при помощи

поиска максимума этой функции.



 Для поиска максимума функции~(\ref{e3-vor}) найдем ее частные

производные по~$a_k$:

 \begin{equation}

 \fr{\partial L(\overline{x},\overline{a})}{\partial a_k} =\sum\limits_{i=1}^M

\fr{\left(dp(x_i\vert k;a_k)/(da_k)\right)\,q_{ij}}

{\sum\limits_{j=1}^N p(x_i\vert j;a_j)q_{ij}}

\,,\enskip k=1,\ldots , N\,.

 \label{e4-vor}

 \end{equation}

Учитывая~(\ref{e2-vor}), получаем производные $p(x\vert k;a_k)$ по~$a_k$:

$$

\fr{dp(x\vert k;a_k)}{a_k} = \fr{1}{\sigma^3\sqrt{2\pi}}\left( x-a_k\right) \exp

\left( -\fr{(x-a_k)^2}{2\sigma^2}\right) \,,\enskip k=1,\ldots , N\,.

$$

Подставляем их в~(\ref{e4-vor}), получаем:

$$

\fr{\partial L(\overline{x},\overline{a})}{\partial a_k} =

\sum\limits_{i=1}^M 

\fr{({1}/{\sigma^2}) \left( x_i-a_k\right) \exp \left( - {(x_i-a_k)^2}/({2\sigma^2})\right)

q_{ik} }{ \sum\limits_{j=1}^N \exp \left( -{(x_i-a_j)^2}/({2\sigma^2})

\right) q_{ij}}\,,\enskip k=1,\ldots , N\,.

$$

Обозначим

$$

c_{ik}(\overline{a}) = \fr{\exp \left( - {(x_i-a_k)^2}/{(2\sigma^2)} \right) q_{ik} }

{ \sum\limits_{j=1}^N \exp \left( - {(x_i-a_k)^2}/{(2\sigma^2)} \right) q_{ij}}

\,.

$$

Тогда

$$

\fr{\partial L(\overline{x},\overline{a})}{\partial a_k}= \sum\limits_{i=1}^M

\fr{1}{\sigma^2} \left( x_i-a_k\right) c_{ik}(\overline{a})\,,\enskip k=1,\ldots , N\,.

$$

Приравняв к нулю производные, получаем систему

\begin{equation}

\sum\limits_{i=1}^M \fr{1}{\sigma^2} \left( x_i-a_k\right)

c_{ik}(\overline{a})=0\,,\enskip k=1,\ldots , N\,.

\label{e5-vor}

\end{equation}

 Выразим из~(\ref{e5-vor}) $a_k$:



\noindent

 $$

 a_k= \sum\limits_{i=1}^M c_{ik}(\overline{a}) x_i \bigg /\sum\limits_{i=1}^M

c_{ik}(\overline{a})\,, \enskip k=1,\ldots , N\,.

 $$

 Введем обозначения:



\noindent

$$ \gamma_{ik}(\overline{a}) = {c_{ik}(\overline{a})}\bigg /{\sum\limits_{i=1}^M

c_{ik}(\overline{a})}\,;

$$



\pagebreak



\noindent

 \begin{align*}

 f_k(\overline{a}) &=\sum\limits_{i=1}^M \gamma_{ik}(\overline{a})x_i\,,\enskip

k=1,\ldots , N\,;\\

 f(\overline{a}) &= (f_1(\overline{a}),\ldots , f_N(\overline{a}))\,.

 \end{align*}

 Тогда система уравнений~(\ref{e5-vor}) примет наиболее простой вид:

 \begin{equation}

 \overline{a}=f(\overline{a})\,.

 \label{e6-vor}

 \end{equation}



 \medskip



 \noindent

 \textbf{Утверждение.}\ \textit{Уравнение~$(\ref{e6-vor})$ имеет решение.}

 \medskip



 \noindent

 Д\,о\,к\,а\,з\,а\,т\,е\,л\,ь\,с\,т\,в\,о\,.\ \ Обозначим

 $$

 Y=\left\{ \overline{y}=(y_1,\ldots , y_N)\vert \min\limits_i x_i\leq

y_l \leq \max\limits_i x_i\,,\ l=1,\ldots , N\right\}\,.

 $$



 Очевидно, что $Y$ является выпуклым и компактным множеством, а

отображение~$f$~--- непрерывным. Так как $0\hm \leq

\gamma_{ik}(\overline{a})\hm\leq 1$, то $\min\limits_i \{x_i\} \hm\leq

f_k(\overline{a})\hm\leq \max\limits_i \{x_i\}$; следовательно, $f:\ Y\hm\to Y$ и

по теореме Брауэра отображение~$f$ имеет неподвижную точку. Утверждение

доказано.

 \smallskip



 Опишем итерационный алгоритм решения уравнения $\overline{a}\hm=

f(\overline{a})$. Пусть задана точность вычислений~$\varepsilon$.



 \smallskip



 \textbf{Алгоритм 1}



 \textit{Положить начальные значения} $\overline{a}^{(0)} =(0, \ldots , 0)$.



 \textit{Вычислять последовательность

$\overline{a}^{(t)}$ по формуле}

 $$

 \overline{a}^{(t+1)} =f(\overline{a}^{(t)})\,,\enskip t=1,2,\ldots,

 $$



 \textit{до тех пор, пока не будет выполнено условие}

$\| \overline{a}^{(t)}-\overline{a}^{(t+1)} \| \hm<\varepsilon$.



 \smallskip



 Рассмотрим обобщение задачи. Пусть имеется несколько базовых станций,

для которых построены карты уровня сигналов. Добавляется несколько новых

базовых станций. Необходимо для них построить карты зон.



 Пусть $d_1,\ldots , d_K$~--- базовые станции, для которых имеются карты

зон, $d_{K+1},\ldots , D_{K+R}$~--- базовые станции, для которых требуется

построить карты зон.



 Предлагается следующий алгоритм решения поставленной задачи.



 \smallskip



 \textbf{Алгоритм 2}



 \textit{Для $t=1,\ldots , R$ выполнить}:



 \hangindent=10mm\hangafter=0\noindent\textit{Получить $M$ наборов измерений

уровня сигнала от базовых станций} $d_1,\ldots , d_{K+t}$.



\hangindent=10mm\hangafter=0\noindent\textit{Используя карты уровня сигналов

для базовых станций $d_1,\ldots , d_{K+t-1}$, вычислить $q_{ij}$,

$i\hm= 1,\ldots , M$, $j\hm= 1,\ldots , N$.}



\hangindent=10mm\hangafter=0\noindent\textit{Используя алгоритм~$1$, построить

карту уровня сигнала для базовой станции}~$d_{K+t}$.



\section{Тестовые испытания}



 Опишем эксперимент, проведенный на реальных данных. В~эксперименте

использовалось оборудование, описанное в табл.~1.



 \begin{table}[h]\small

 \begin{center}

 \Caption{Использованное оборудование}

 \vspace*{2ex}



 \tabcolsep=3pt

 \begin{tabular}{|l|l|}

 \hline

 \multicolumn{1}{|c|}{Оборудование}&\multicolumn{1}{c|}{Описание}\\

 \hline

\tabcolsep=0pt\begin{tabular}{l}Базовая\\ станция \\(4~шт.)\end{tabular}&\tabcolsep=0pt\begin{tabular}{l}Беспроводная точка доступа Netgear WAG 102\\

Мощность передатчика 100 мВт\\

1 внешняя антенна 5 dBi (2,4 ГГц)\end{tabular}\\

\hline

Ноутбук&\tabcolsep=0pt\begin{tabular}{l}MacBook Pro, Core i5 2,4~ГГц, 4ГБ ОЗУ, WiFi

802.11n (Broadcom)\\

Mac OS X 10.7.5, 64-битная архитектура\end{tabular}\\

\hline

\end{tabular}

\end{center}

\end{table}



 На территории физико-тех\-ни\-че\-ско\-го факультета Петрозаводского

государственного университета была выбрана тестовая площадка размером

$12\times20$~м, на которой было задействовано 4~базовые станции ($d_1, \ldots

, d_4$, $D\hm=4$, рис.~3).

 Данная площадка имеет опорные колонны, что

позволяет эмитировать условия непрямой видимости какой-либо из базовых

станций мобильным объектом. Для\linebreak

\vspace*{-12pt}



 \begin{floatingfigure}{65mm} %fig3

 \vspace*{1pt}

 \begin{center}

 \mbox{%

 \epsfxsize=60.198mm

 \epsfbox{vor-3.eps}

 }

 \end{center}

 \vspace*{-9pt}

 \Caption{Опорные точки тестовой площадки}

 \vspace*{9pt}

 \end{floatingfigure}



\noindent

определения оценок математических

ожиданий~$a_{jd}$ уровней сигналов базовых станций были проведены

измерения в 22~опорных точках ($N\hm=22$), по 50~замеров в четырех

направлениях относительно стен здания (итого 200~замеров в каждой опорной

точке).



Полученные в ходе проведения измерений карты уровня сигналов для базовых

станций $d_1,\ldots , d_4$ пред\-став\-ле\-ны на рис.~4. Вертикальная ось

соответствует уровню сигнала (дБм), горизонтальная плоскость соответствует

размерам тестовой площадки, маркерами отмечены реальные усредненные

значения.



\begin{figure} %fig4

 \vspace*{1pt}

 \begin{center}

 \mbox{%

 \epsfxsize=123.767mm

 \epsfbox{vor-4.eps}

 }

 \end{center}

 \vspace*{-9pt}

\Caption{Эмпирические карты уровня сигналов:

(\textit{а})~$d_1$;

(\textit{б})~$d_2$;

(\textit{в})~$d_3$;

(\textit{г})~$d_4$

}

%\end{figure}

\vspace*{10pt}

%\begin{figure} %fig5

 \vspace*{1pt}

 \begin{center}

 \mbox{%

 \epsfxsize=124.325mm

 \epsfbox{vor-5.eps}

 }

 \end{center}

 \vspace*{-9pt}

\Caption{Смоделированные карты уровня сигналов:

(\textit{а})~$d_1$;

(\textit{б})~$d_2$;

(\textit{в})~$d_3$;

(\textit{г})~$d_4$

}

\end{figure}





 Далее были использованы те же измерения для тестирования работы

алгоритма~1 ($M\hm= 4400$). Используя данные от трех из четырех базовых

станций, по формуле~(\ref{e1-vor}) в момент каждого измерения~$i$

вычислялись вероятности~$q_{ij}$ на\-хож\-де\-ния объекта в опорных точках

$j\hm= 1,\ldots , N$. Затем с помощью алгоритма~1 вычислялись оценки

математических ожиданий~$a_j$ в опорных точках $j\hm= 1,\ldots , N$ для

четвертой базовой станции. Параметр~$\sigma$ был равен трем.



 Таким образом, были смоделированы карты уровня сигналов для каждой

базовой станции $d_1,\ldots , d_4$. Пример смоделированных карт уровня

сигналов для всех базовых станций представлен на рис.~5. Найденные

значения~$a_j$ сравнивались с реальными средними значениями уровня

сигнала~$a_j^*$ для данной базовой станции. Полученное в ходе эксперимента

число опорных точек с различным значением $\vert a_j\hm- a_j^*\vert$

представлено в табл.~2.



\begin{table}\small

\begin{center}



\parbox{275pt}{\Caption{Разница между реальными и спрогнозированными математическими ожиданиями

уровня сигнала}



}

\vspace*{2ex}



\begin{tabular}{|c|c|c|}

\hline

&&\\[-9pt]

$\vert a_j-a_j^*\vert$&Число опорных точек&

Доля опорных точек, \%\\[3pt]

\hline

$[$0--1)&35&39,8\\

$[$1--2)&24&27,3\\

$[$2--3)&11&12,5\\

$[$3--4)&11&12,5\\

4 и больше&\hphantom{9}7&\hphantom{9}7,9\\

\hline

\end{tabular}

\end{center}

\end{table}





 В двух случаях был зафиксирован выброс $\vert a_j\hm- a_j^*\vert \hm\approx

16$. Однако следует отметить, что такие большие отклонения возможны и при

<<ручном>> способе составления карты уровня сигналов.



 Для определения сходимости работы алгоритма~2 был проведен

имитационный эксперимент, заключающийся в следующем. Для десяти

базовых станций ($D\hm = 10$) были назначены условные математические

ожидания уровня сигналов в десяти опорных точках ($N\hm = 10$).

Параметр~$\sigma$ был равен трем. В~ходе каждого эксперимента выбиралось

значение $K\hm<10$, на основе карт уровня сигналов первых $K$ базовых

станций строились карты уровня сигналов $R$ остальных базовых станций

($R\hm= D \hm- K)$. В~каждом из $M$~измерений ($M\hm = 1000$) случайным

образом выбиралась опорная точка и разыгрывались нормально

распределенные случайные величины~$x_{id}$~--- уровни сигнала от базовых

станций $d\hm=1,\ldots , D$ в данной опорной точке. Уровни сигналов~$x_{id}$

поступали на вход алгоритма~2, при помощи которого на основе карт уровня

сигналов первых $K$ базовых станций определялись карты уровня сигналов

для базовых станций $d\hm= K\hm+1,\ldots , D$.



 Далее полученные результаты использовались для определения позиции

мобильного объекта следующим образом. Случайным образом выбиралась

опорная точка, генерировались сигналы и определялось местоположение точки

при помощи карт уровня сигналов всех базовых станций (начальных и

обученных). Для этого по формуле~(\ref{e1-vor}) вычислялись апостериорные

вероятности опорных точек и в качестве искомой определялась точка с

максимальной вероятностью. После этого проверялось, совпала ли найденная

таким образом опорная точка с заранее выбранной. Доля неправильно

определенных опорных точек в ходе всего имитационного эксперимента

представлена в табл.~3.



 \begin{table}\small

 \begin{center}



\parbox{265pt}{\Caption{Доля неправильно определенных опорных точек при различном числе

начальных базовых станций}



}

 \vspace*{2ex}



 \begin{tabular}{|c|c|c|c|c|}

 \hline

&\multicolumn{4}{c|}{\ }\\[-9pt]

Число добавленных

&\multicolumn{4}{c|}{Начальное число базовых станций~$K$}\\

\cline{2-5}

&&&&\\[-9pt]

базовых станций $R$

&9&7&5&3\\

\hline

&&&&\\[-10pt]

1&0,088&0,121&0,204&0,285\\

2& &0,13\hphantom{9}&0,17\hphantom{9}&0,227\\

3& &0,114&0,148&0,213\\

4& & &0,128&0,163\\

5& & &0,077&0,144\\

6& & & &0,124\\

7& & & &0,105\\

\hline

\end{tabular}

\end{center}

\end{table}



 Из табл.~3 видно, что точность определения опорных точек улучшается при

добавлении новых базовых станций, для которых карты уровня принимаемого

сигнала строятся при помощи алгоритма~2.



 \section{Заключение}



 Результаты полученных расчетов показали, что при уменьшении с 9 до 3

числа исходных <<обученных>> базовых станций число неправильно

определенных опорных точек увеличилось менее чем на 2\%. Это указывает на

практическую применимость описанной модели для формирования карт уровня

сигналов в ходе эксплуатации системы и их дальнейшего использования при

уточнении позиций мобильных объектов.



 Полученные в данной работе результаты также могут быть применены для

построения систем локации с резервными базовыми станциями, для которых

отсутствуют карты уровня сигналов. Кроме того, описанный алгоритм можно

использовать для корректировки уже существующих карт уровня сигналов.



 В дальнейшем планируется исследовать задачу построения карт уровня

сигналов устройств различных производителей на основе имеющейся карты

уровня сигналов одного из этих устройств.



{\small\frenchspacing

{%\baselineskip=10.8pt

\addcontentsline{toc}{section}{Литература}

\begin{thebibliography}{9}

\bibitem{2-vor}

\Au{Galov A.\,S., Moschevikin A.\,P., Voronov~R.} Combination of RSS

localization and ToF ranging for increasing positioning accuracy indoors~/\!\!/

ITST-2011: 11th Conference (International) on ITS Telecommunications

Proceedings.~--- St. Petersburg: IEEE, 2011. P.~299--304.

\bibitem{1-vor}

\Au{Moschevikin A., Voronov R., Galov~A., Soloviev~A.} Using pressure sensors

for floor identification in wireless sensors networks~/\!\!/ IDAACS-SWS: 1st

Symposium (International) on Wireless Systems Proceedings.

%Offenburg, German. September 20--21, 2012.~

--- Offenburg: IEEE, 2012. P.~2--6.

\bibitem{3-vor}

\Au{Galov A., Moschevikin A.} Bayesian filters for ToF and RSS measurements

for indoor positioning of a mobile object~/\!\!/ IPIN-2013: Conference

(International) on Indoor Positioning and Indoor Navigation Proceedings.

%Montbeliard, France. October 28--31, 2013.~

--- Montbeliard: IEEE, 2013.

P.~310--317.

\bibitem{4-vor}

\Au{Under J.} Deploying license-free wireless wide-area networks.~---

Indianapolis, Indiana: Cisco Press, 2004. 333~p.

\bibitem{5-vor}

\Au{Tomsho G.} Guide to networking essentials.~--- 6th ed.~--- Stamford,

Connecticut: Cengage Learning, 2011. 752~p.

\bibitem{6-vor}

\Au{Hassan A.} Advanced location-based technologies and services.~--- Boca

Raton, Florida: CRC Press, 2013. 354~p.



\end{thebibliography}

}

}



\hfill{\small\textit{Поступила в редакцию 21.12.13}}



\vspace*{14pt}



\hrule



\vspace*{2pt}



\hrule

\def\tit{DYNAMIC CREATION OF WiFi-SIGNAL LEVEL MAPS FOR~LOCATION

SYSTEMS}



\def\titkol{Dynamic creation of WiFi-signal level maps for~location

systems}



\def\aut{R.\,V.~Voronov and S.\,V.~Malodushev}

\def\autkol{R.\,V.~Voronov and S.\,V.~Malodushev}





\titel{\tit}{\aut}{\autkol}{\titkol}



\vspace*{-15pt}



\noindent

\begin{center}

Petrozavodsk State University, 33 Lenin Prosp.,

Petrozavodsk 185910, Karelia,

Russian Federation

\end{center}





\def\leftfootline{\small{\textbf{\thepage}

\hfill Sistemy i Sredstva Informatiki~---

Systems and Means of Informatics 2014 vol~24 no\,1}

}%

 \def\rightfootline{\small{Sistemy i Sredstva Informatiki~---

 Systems and Means of Informatics 2014 vol~24 no\,1

\hfill \textbf{\thepage}}}



\vspace*{3pt}



\Abste{Signal level maps are designed to locate wireless mobile devices. To increase

a wireless network segment, it is necessary to create new signal level maps for

implementation of location of wireless mobile devices in new areas. Traditionally,

such maps are based on a great number of measurements taken at certain

points, the process of these measurements requires a lot of time. In the article,

dynamic creation of signal level maps is investigated. A~heuristic

algorithm for dynamic creation of signal level maps during operation of a location

system is described. A~function similar to the likelihood function was used for estimation

of the average values of the signal levels. WiFi networks

operating in the 2.4-gigahertz frequency band were used as wireless networks.

 Test results show that the suggested

algorithm can be used for solving this issue.}



\vspace*{1pt}



\KWE{WiFi; base station; control point; signal level; location; signal level map}



\vspace*{1pt}



\DOI{10.14357\!/08696527140105}





\Ack

\noindent

The research was done in the frames of MIP Nanonets Ltd.\ and RTL-Service Co.\ activity.

The research was also supported by the Petrozavodsk State University (PSU)

(Program of strategy development of PSU in 2012--2016), Russian Ministry of Education

and Science, SR MFP Foundation and NTS and Ministry of Economic Development of Karelia 

(Program ``Development of sky system for mobile abonents of distributed WiFi network'' in

2013) as well as in the frames of program UMNIK (Member of the Youth Research and

Innovation Competition) ``Analysis and design

of WiFi networks inside buildings'' in 2012--2013.





\renewcommand{\bibname}{\protect\rmfamily References}



{\small\frenchspacing

{%\baselineskip=10.8pt

%\addcontentsline{toc}{section}{Литература}

\begin{thebibliography}{9}

\bibitem{1-vor-1}

\Aue{Moschevikin, A., R.~Voronov, A.~Galov, and A.~Soloviev}. 2012.

Using pressure sensors for floor identification in wireless sensors networks.

\textit{1st Symposium (International) on Wireless Systems IDAACS-SWS Proceedings}.

Offenburg, Germany: IEEE. 2--6.

\bibitem{2-vor-1}

\Aue{Galov, A.\,S., A.\,P.~Moschevikin, and R.~Voronov}. 2011.

Combination of RSS localization and ToF ranging for increasing positioning accuracy indoors.

\textit{11th Conference (International) on ITS Telecommunications ITST-2011 Proceedings}.

St.\ Petersburg: IEEE. 299--304.

\bibitem{3-vor-1}

\Aue{Galov, A., and A.~Moschevikin}. 2013. Bayesian filters for ToF and RSS

measurements for indoor positioning of a mobile object.

\textit{Conference (International) on Indoor Positioning and Indoor Navigation (IPIN-2013)

Proceedings}. Montbeliard, France: IEEE. 310--317.

\bibitem{4-vor-1}

\Aue{Under, J.} 2004. \textit{Deploying license-free wireless wide-area networks}.

Indianopolis, Indiana:  Cisco Press.  333~p. 

\bibitem{5-vor-1}

\Aue{Tomsho, G.} 2011. Guide to networking essentials. 6th ed.

Stamford, Connecticut: Cengage Learning. 752~p. 

\bibitem{6-vor-1}

\Aue{Hassan, A.} 2013.

\textit{Advanced location-based technologies and services}.

Boca Raton, Florida: CRC Press. 354~p. 

 \end{thebibliography}

} }



\vspace*{-6pt}



\hfill{\small\textit{Received December 21, 2013}}



\vspace*{-6pt}



\Contr



\noindent

\textbf{Voronov Roman V.} (b.\ 1976)~--- Candidate of Science (PhD) in technology,

associate professor, Petrozavodsk State University, 33 Lenin Prosp.,

Petrozavodsk 185910, Karelia,

Russian Federation; rvoronov@sampo.ru



\vspace*{3pt}



\noindent

\textbf{Malodushev Sergei V.} (b.\ 1985)~--- leading engineer,

Petrozavodsk State University, 33 Lenin Prosp.,

Petrozavodsk 185910, Karelia,

Russian Federation; svm@petrsu.ru







\label{end\stat}



\renewcommand{\bibname}{\protect\rmfamily Литература}

